\section{Option B --- Alfred Discovery Architecture (Overview)}
\label{sec:option-b-overview}

\subsection{System shape}

\begin{diagram}
                     +----------------------+
                     |   ALFRED ANALYZER    |
                     |  (aggregation +      |
                     |   discovery engine)  |
                     +-----------+----------+
                                 |
                        Gigabit Ethernet (host-side,
                        isolated from vehicle network)
                                 |
        +------------------+----+----+------------------+
        |                  |         |                  |
   +----+-----+      +-----+----+ +--+-------+     +-----+-----+
   |Probe Node|      |Probe Node| |Probe Node|     |Probe Node |
   |  A: CAN/ |      |  B: Auto |     |  C: LIN/ |     |  D: Analog/|
   |  CAN-FD  |      |  Ethernet|     |  UART    |     |  PWM/Hall  |
   +----+-----+      +-----+----+ +--+-------+     +-----+-----+
        |                  |         |                  |
   OBD-II / splice     TAP on         splice into      probe points on
   into CAN-B, CAN-C   100/1000BASE-T1  LIN-2, LIN-3    motor leads,
   harness (dev unit)  segment (dev)   body network      Hall sensor test
                                                          points
\end{diagram}

The probe is explicitly \emph{distributed}, not a single OBD-II dongle, because OBD-II is a single point behind at least one gateway and typically exposes only diagnostic-relevant traffic on one or two buses. A distributed set of probe nodes, each timestamp-synchronized to the others, can observe multiple physically separate buses (CAN-B body network, CAN-C powertrain, a LIN sub-bus behind a body controller, an automotive Ethernet backbone segment) simultaneously and correlate across them --- which is the entire basis of the semantic-identification workflow in Section~\ref{sec:probe-discovery}.

\subsection{The passive/active boundary}

\begin{decision}[title={Non-negotiable design boundary}]
Passive capture and active experimentation are architecturally separate modes with separate hardware enable paths (Section~\ref{sec:active-control}). A probe node cannot transmit unless: (1) it is explicitly placed in Active Test Mode by an operator, (2) the target signal has a validated mapping with confidence above threshold, (3) the command passes an allowlist and state-validation policy, and (4) a hardware transmit-enable line (not just a software flag) is asserted. Default state on power-up is receive-only.
\end{decision}

\subsection{From bytes to behavior}

Option B is a four-stage pipeline, detailed in Sections~\ref{sec:probe-discovery}--\ref{sec:knowledge-graph}: electrical/protocol identification $\rightarrow$ network protocol recognition $\rightarrow$ semantic identification (via structured experiments and/or imported engineering files) $\rightarrow$ Vehicle Knowledge Graph, which is then compiled into the same Machine IR that Option A populates directly from E2B capability manifests. That equivalence is what makes convergence (Section~\ref{sec:convergence}) real rather than aspirational.
